\documentclass[11pt]{article}

\usepackage[margin=1in]{geometry}
\usepackage{amsmath}
\usepackage{booktabs}
\usepackage{enumitem}
\usepackage{graphicx}
\usepackage{siunitx}
\usepackage{hyperref}
\usepackage{gensymb}



\title{Intermediate Crawling Load Test}
\author{Rift / Team 09}
\date{\today}

\begin{document}


\maketitle

\section*{Test ID}
\textbf{T-01ICa Intermediate Crawling Test}

\section*{Test Objective}
Quantify the maximum weight the payload can handle while performing the test criteria of the Basic Crawling Test (T-01BC). 

\section*{Robot Configuration}
\begin{itemize}[noitemsep]
    \item Robot configurations: Standing
    \item Tube pressures: 15 psi
    \item Control mode: Open-loop
    \item Electronics status: Batteries are fully charged (Greater than 4V per cell, greater than 12 Volts per battery)
\end{itemize}

\section*{Materials}
\begin{itemize}[noitemsep]
    \item Plywood payload container (mass = ...)
    \item Weights (incrementing by 1 kg, 5 kg, 10 kg)
    \item Stopwatch timer
    \item Measuring tape (at least 5 m in length)
    
\end{itemize}

\section*{Environment and Setup}
\begin{itemize}[noitemsep]
    \item Surface type: Smooth Finished Concrete
    \item Surface angle: $\psi = 0 {\degree} $
    \item External load: Variable
    \item Measurement method: time the rover moving directly and laterally through 2 full step cycles under different loads. A video of the test will also be recorded for future analysis.
\end{itemize}

\section*{Initial Conditions}
\begin{itemize}[noitemsep]
    \item Initial position: $x_0 = 0$
    \item Initial orientation: 
    \begin{itemize}
        \item Direct crawling: $\theta_z = 0 {\degree}$
        \item Lateral crawling:  $\theta_z = 90 {\degree}$
    \end{itemize}
    \item Initial shape state: Standing
    \item Stabilization time before command: 10?~s
\end{itemize}

\section*{Command Sequence}
\begin{itemize}[noitemsep]
    \item Command type: \underline{\hspace{5cm}}
    \item Command values: \underline{\hspace{5cm}}
    \item Update rate: 3~Hz
    \item Command Steps: 8~s
\end{itemize}

\section*{Measured Quantities}
\begin{itemize}[noitemsep]
    \item Mass contained in the payload before rover cannot move on flat ground
    \item Average node height during direct and lateral traversal under load
    \item Speed of the rover under different loading conditions
\end{itemize}

\section*{Success Criteria (Per Trial)}
\begin{itemize}[noitemsep]
    \item Forward traversal speed: .1 km/h  $\leq$ .14 km/h $\leq$ No upper limit
    \item Side to Side traversal speed: ?? km/h  $\leq$ ?? km/h $\leq$ No upper limit
    \item  Step height: 10 cm  $\leq$ 25 cm $\leq$ No upper limit
\end{itemize}

\section*{Test Procedure}
\begin{enumerate}[noitemsep]
    \item Place robot in standing position on an ungraded surface
    \item Complete pretest checklist
    \item Place maximum mass found from the Basic Static Test (T-01BS) into the payload
    \item Send command sequence to drive rover 8 steps (2 full step cycles) in the positive x-direction
    \item Stop timer upon completion of sequence
    \item Measure distance between front left and front right nodes before and after step sequence.
    \item Record average of the two above measured distances divided by traversal time
    \item Reposition rover to the start of path
    \item Repeat steps 1-7, driving the rover laterally instead of directly
    \item If the rover fails to move while handling the current test load, increment the mass back in increments of 10 kg, 5 kg, and then 1 kg. Repeat steps 1-9 with adjusted loads. With the maximum load found that the rover can handle while crawling, move onto steps 11+.
    \item Reset robot to standing position with maximum permitted crawling load
    \item Send command sequence to drive rover halfway through step
    \item Measure height of stepping node
    \item Repeat steps 9 and 10 through 8 sets of steps
    \item Record average of step heights
    \item Repeat steps 10-14 with lateral travel under the maximum lateral travel load

\end{enumerate}

\section*{Number of Trials}
\begin{itemize}[noitemsep]
    \item Total trials: 1
    \item Time between trials: N/A
\end{itemize}

\section*{Results}
Summarize the observed performance metrics.

\begin{center}
\begin{tabular}{llc}
 & Direct Crawling ($\dot x$)&Lateral Crawling ($\dot y$)\\
\midrule
Traversal Speed (km/h)&   &\\
 Maximum Load (kg)& &\\
 Average Step Height (cm)& &\\
\end{tabular}
\end{center}

\section*{Observed Failure Modes}
Describe any observed failure behaviors.

\begin{itemize}[noitemsep]
    \item Inflated beams of the truss buckle during stepping
    \item Motors stall of fail to provide sufficient torque to move the truss
    \item Truss fails to lift nodes from the ground
    \item Truss tips prematurely while lifting node from the ground
\end{itemize}

\section*{Conclusions}
Provide a concise interpretation of the results.

\textit{Example:} The robot achieved consistent forward motion on flat ground with an average speed of \underline{\hspace{1cm}}~cm/s while under a load of \underline{\hspace{1cm}}~cm. This cooresponds to a \underline{\hspace{1cm}} \% reduction from the speed with no load. 

\section*{Notes and Limitations}
Document sources of variability or limitations.

\end{document}